\section{Representation Learning as a Unifying Principle}

\subsection{Overview}

The preceding sections introduced representation as a map $\phi: \mathcal{X} \to \mathcal{Z}$, self-supervised objectives that shape $\phi$, and transfer mechanisms that reuse it.  
This section elevates these ideas into a unifying principle:

\medskip

\noindent
\textbf{Thesis.}  
Modern machine learning can be understood as the construction of intermediate coordinate systems $\mathcal{Z}$ in which multiple tasks become simpler, more stable, or more transferable.

\medskip

This section formalizes key properties of representations and shows how they operate across supervised learning, self-supervised learning, generative modeling, and reinforcement learning.

\subsection{Core Properties of Representations}

\begin{definition}[Invariance]
A representation $\phi$ is invariant to a transformation $T:\mathcal{X}\to\mathcal{X}$ if
\[
\phi(Tx) = \phi(x) \quad \text{for all } x \in \mathcal{X}.
\]
\end{definition}

\begin{definition}[Sufficiency]
A representation $\phi$ is sufficient for a task $y$ if there exists a predictor $g$ such that
\[
y = g(\phi(x)),
\]
or more generally, the conditional distribution satisfies
\[
P(y \mid x) = P(y \mid \phi(x)).
\]
\end{definition}

\begin{definition}[Transferability]
A representation $\phi$ is transferable across a family of tasks $\mathcal{T}$ if for each $T \in \mathcal{T}$ there exists a simple predictor $g_T$ (e.g., linear or low-complexity) such that
\[
f_T(x) = g_T(\phi(x)).
\]
\end{definition}

\begin{definition}[Semantic Structure]
A representation exhibits semantic structure if distances or directions in $\mathcal{Z}$ correspond to meaningful variations in the data-generating process.
\end{definition}

\begin{remark}
These properties are not independent. Increasing invariance may reduce sufficiency, and transferability depends on both.
\end{remark}

\subsection{A Cross-Domain Representation Framework}

Across learning paradigms, we observe a shared decomposition:
\[
x \xmapsto{\;\phi\;} z \xmapsto{\;g\;} y.
\]

\medskip

\noindent
\textbf{Supervised learning.}
\begin{itemize}
\item $\phi$ is learned jointly with $g$
\item sufficiency is enforced by labels
\item invariance emerges implicitly
\end{itemize}

\medskip

\noindent
\textbf{Self-supervised learning.}
\begin{itemize}
\item $\phi$ is trained via predictive objectives
\item invariance is explicitly induced (augmentations, masking)
\item transferability is evaluated downstream
\end{itemize}

\medskip

\noindent
\textbf{Generative modeling.}
\begin{itemize}
\item $\phi$ corresponds to latent variables
\item structure is imposed through likelihood or score objectives
\item sufficiency relates to reconstruction or density estimation
\end{itemize}

\medskip

\noindent
\textbf{Reinforcement learning.}
\begin{itemize}
\item $\phi$ defines state representations
\item sufficiency relates to Markov structure
\item transferability appears in policy reuse and generalization
\end{itemize}

\begin{proposition}[Induced Simplicity Across Tasks]
If $\phi$ is sufficient for a family of tasks $\mathcal{T}$ and invariant to nuisance transformations common to $\mathcal{T}$, then each task admits a predictor $g_T$ of reduced complexity compared to operating directly on $\mathcal{X}$.
\end{proposition}

\begin{proof}[Sketch]
Sufficiency ensures no relevant information is lost.  
Invariance removes irrelevant variation.  
Thus the effective input space for $g_T$ is both smaller and more structured, reducing hypothesis complexity.
\end{proof}

\subsection{Worked Examples}

\paragraph{Example 1: Vision (Invariance)}
Let $x$ be an image and $T$ a small translation.  
A convolutional representation $\phi$ approximately satisfies
\[
\phi(Tx) \approx \phi(x),
\]
making object classification invariant to position.

\paragraph{Example 2: Language (Contextual Representation)}
A token embedding $\phi(w_i)$ depends on context:
\[
\phi(w_i) = f(w_1, \dots, w_T).
\]
Syntactic and semantic relations become approximately linear or geometrically structured in $\mathcal{Z}$.

\paragraph{Example 3: Reinforcement Learning (State Abstraction)}
Let $x$ be raw observations and $\phi(x)$ a learned state.  
If
\[
P(s_{t+1} \mid s_t, a_t) \approx P(s_{t+1} \mid x_t, a_t),
\]
then $\phi$ defines a near-Markov representation, simplifying control.

\subsection{Representation as Geometry}

A useful abstraction is to view $\phi$ as defining a coordinate system on data:

\[
\mathcal{X} \xrightarrow{\;\phi\;} \mathcal{Z},
\]

where:
\begin{itemize}
\item distances reflect similarity,
\item directions reflect variation factors,
\item subspaces correspond to task-relevant structure.
\end{itemize}

Multiple tasks then act as functions on the same latent space:
\[
g_1, g_2, g_3 : \mathcal{Z} \to \mathcal{Y}.
\]

\medskip

\noindent
\textbf{Interpretation.}  
Representation learning constructs a shared geometric substrate on which diverse tasks are defined.

\subsection{Remarks: Limits of the Representation View}

\begin{remark}[Non-identifiability]
Many representations yield identical performance.  
There is no unique ``correct'' latent space.
\end{remark}

\begin{remark}[Interpretability is not guaranteed]
Semantic structure is not automatic.  
A representation may be useful for prediction while remaining opaque.
\end{remark}

\begin{remark}[Task-dependence]
A representation that is optimal for one task may discard information needed for another.
\end{remark}

\begin{remark}[Overclaiming risk]
Statements such as ``the model understands concepts'' should be treated cautiously.  
Most learned representations are aligned with objectives, not necessarily with human semantics.
\end{remark}

\subsection{Synthesis}

This chapter can now be read as a progression:

\begin{itemize}
\item Section 4.1: representations as maps
\item Section 4.2: objectives that shape representations
\item Section 4.3: reuse through transfer
\item Section 4.4: scaling representation learning to foundation models
\item Section 4.5: unification across paradigms
\end{itemize}

\medskip

\noindent
\textbf{Conclusion.}  
Representation learning is not a subfield but a structural principle:  
learning useful coordinate systems in which problems become simpler.

\subsection{What to Remember}

\begin{itemize}
\item Representations balance invariance and sufficiency
\item Transferability is a central design goal
\item Many learning paradigms share the same representation structure
\item Representations define geometry, not just features
\end{itemize}

\subsection{Common Confusion}

\begin{itemize}
\item Invariance does not mean information loss is harmless
\item Good performance does not imply interpretable representations
\item Transferability depends on task families, not universality
\end{itemize}

\subsection{Exercises}

\begin{exercise}
Give an example where increasing invariance reduces performance due to loss of sufficiency.
\end{exercise}

\begin{exercise}
Show that if $\phi$ is invertible, it is always sufficient but not necessarily invariant.
\end{exercise}

\begin{exercise}
Compare representations learned by supervised and contrastive objectives on the same dataset.
\end{exercise}

\begin{exercise}
Construct a simple RL example where a poor representation violates the Markov property.
\end{exercise}

\subsection{Notes and References}

Key directions include:
\begin{itemize}
\item early representation learning and feature learning literature
\item invariance and equivariance in convolutional architectures
\item contrastive and self-supervised learning frameworks
\item latent-variable models in generative modeling
\item representation learning in reinforcement learning
\end{itemize}

This area remains partially formalized, with strong empirical guidance and evolving theoretical foundations.